
طیف سیگنال داده شده است:

\[
X(j\omega)=
\begin{cases}
1-\dfrac{|\omega|}{200\pi}, & |\omega|\le 200\pi\\[2mm]
0, & |\omega|>200\pi
\end{cases}
\]

که یک طیف مثلثی با پهنای باند
\(
\omega_m=200\pi
\)
است.


\subsectionaddtolist{الف}


\[
p(t)=\sum_{n=-\infty}^{\infty}\delta(t-nT)
\]

طبق زوج تبدیل فوریه قطار ضربه:

\[
P(j\omega)
=
\frac{2\pi}{T}
\sum_{k=-\infty}^{\infty}
\delta(\omega-k\omega_s)
\]

\[
{
P(j\omega)=
\frac{2\pi}{T}
\sum_{k=-\infty}^{\infty}
\delta(\omega-k\omega_s)
}
\]

%%%%%%%%%%%%%%%%%%%%%%%%%%%%%%%%%%%%%%%%%%%%%%%%%%%%%%

\subsectionaddtolist{ب}

چون
\(
x_p(t)=x(t)p(t)
\)
استفاده از خاصیت کانولوشن در حوزه فرکانس:
\[
X_p(j\omega) = \frac{1}{2\pi} X(j\omega) * P(j\omega)
\]

جایگذاری \(P(j\omega)\) داریم:

\[
X_p(j\omega)
=
\frac{1}{2\pi}
X(j\omega)
*
\left(
\frac{2\pi}{T}
\sum_{k=-\infty}^{\infty}
\delta(\omega-k\omega_s)
\right)
\]

بنابراین:

\[
{
X_p(j\omega)
=
\frac{1}{T}
\sum_{k=-\infty}^{\infty}
X\!\left(j(\omega-k\omega_s)\right)
}
\]

یعنی طیف اولیه با فاصله‌های $\omega_s$ تکرار شده و ضریب
$\frac1T$
گرفته است.

%%%%%%%%%%%%%%%%%%%%%%%%%%%%%%%%%%%%%%%%%%%%%%%%%%%%%%

\subsectionaddtolist{ج}

ابتدا فرکانس نمونه‌برداری:

\[
\omega_s
=
\frac{2\pi}{T}
=
600\pi
\]

از آنجا که
\(
\omega_m=200\pi,
\)
داریم:

\[
\omega_s=600\pi >2\omega_m=400\pi
\]

پس هیچ همپوشانی (Aliasing) رخ نمی‌دهد.

در خروجی فیلتر پایین‌گذر ایده‌آل با
\(
\omega_c=300\pi,
\)
فقط نسخه مرکزی طیف عبور می‌کند، لذا:

\[
Y(j\omega)
=
\frac1T X(j\omega)
\]

بنابراین در حوزه زمان:

\[
{
y(t)=\frac1T\,x(t)
}
\]

و چون
\(
T=\frac1{300},
\)
نتیجه
\(
y(t)=300\,x(t)
\)
است.

%%%%%%%%%%%%%%%%%%%%%%%%%%%%%%%%%%%%%%%%%%%%%%%%%%%%%%

\subsectionaddtolist{د}
\[x(t)=\cos(150\pi t)+\cos(300\pi t)\]

اکنون:

\[
\omega_s
=
\frac{2\pi}{T}
=
200\pi
\]

طیف ورودی:

\[
X(j\omega)
=
\pi\Big[
\delta(\omega-150\pi)+\delta(\omega+150\pi)
\Big]
+
\pi\Big[
\delta(\omega-300\pi)+\delta(\omega+300\pi)
\Big]
\]

پس از نمونه‌برداری:

\[
X_p(j\omega)
=
100
\sum_{k=-\infty}^{\infty}
X\!\left(j(\omega-k200\pi)\right).
\]

برای مؤلفه
\(
\omega=150\pi
\)
نسخه‌های طیفی در
\(
\pm150\pi,\;
\pm50\pi,\;
\pm350\pi,\ldots
\)
ظاهر می‌شوند.

برای مؤلفه
\(
\omega=300\pi
\)
نسخه‌ها در
\(
\pm300\pi,\;
\pm100\pi,\;
\pm500\pi,\ldots
\)
قرار می‌گیرند.

فیلتر پایین‌گذر با
\(
\omega_c=300\pi
\)
فقط مؤلفه‌های داخل بازه
\(
|\omega|\le 300\pi
\)
را عبور می‌دهد.

در نتیجه در خروجی فرکانس‌های زیر وجود دارند:

\[
50\pi,\qquad
100\pi,\qquad
150\pi,\qquad
300\pi.
\]

و چون هر کپی طیفی ضریب
\(
\frac1T=100
\)
دارد، خروجی برابر است با:

\[
{
y(t)
=
100\Big[
\cos(50\pi t)
+
\cos(100\pi t)
+
\cos(150\pi t)
+
\cos(300\pi t)
\Big]
}
\]

\subsectionaddtolist{توضیح پدیده نامطلوب}

چون
\(
\omega_s=200\pi
<
2(300\pi)
\)
شرط نایکوئیست نقض شده است و نسخه‌های تکراری طیف روی یکدیگر می‌افتند.

در نتیجه مؤلفه
\(
300\pi
\)
به فرکانس
\(
|300\pi-200\pi|
=
100\pi
\)
تاخورده می‌شود و مؤلفه
\(
150\pi
\)
نیز تصویری در
\(
50\pi
\)
ایجاد می‌کند.

این پدیده همان
{\text{Aliasing (همپوشانی طیفی)}}
است که باعث می‌شود سیگنال اصلی از روی نمونه‌ها به طور یکتا قابل بازیابی نباشد.
